\documentclass[../DoAn.tex]{subfiles}

\begin{document}

\section{Tổng quan chương}

Chương này trình bày các nền tảng lý thuyết và công nghệ được sử dụng trực tiếp
trong đồ án. Nội dung không nhằm khảo sát toàn bộ lĩnh vực xử lý PDF hay hỏi đáp
tài liệu, mà chỉ tập trung vào các kỹ thuật xuất hiện trong quy trình đã thiết kế và
thực nghiệm.

Nội dung chương được chia theo cấu trúc hiện tại của đồ án. Nhóm thứ nhất làm
cơ sở cho Chương 4, gồm xử lý PDF, phân tích bố cục, trích xuất bảng, chia đoạn
và lập chỉ mục. Nhóm thứ hai làm cơ sở cho Chương 5, gồm truy xuất thông tin,
hỏi đáp có căn cứ và dẫn chứng nguồn. Phần cuối trình bày các độ đo được dùng
trong Chương 6 để đánh giá từng tầng của hệ thống.

\section{Xử lý PDF và phân tích bố cục tài liệu}
\label{section:pdf-layout-foundation}

PDF ưu tiên bảo toàn hình thức hiển thị, nên nội dung thường gắn với tọa độ hình
học trên trang. Vì vậy, khi xử lý PDF cho truy xuất và hỏi đáp, hệ thống cần chuyển
nội dung từ dạng hiển thị sang các đơn vị có cấu trúc như trang, vùng nội dung,
block và đơn vị truy xuất. Các bộ dữ liệu như DocLayNet và PubLayNet cho thấy nhận diện
bố cục tài liệu là một bài toán quan trọng trong khai thác PDF
\cite{pfitzmann2022doclaynet,zhong2019publaynet}.

Trong đồ án, mỗi vùng nội dung được biểu diễn bằng khung bao
(\textit{bounding box}), loại nội dung và siêu dữ liệu nguồn. Với PDF có lớp văn bản,
hệ thống lấy văn bản và tọa độ bằng thư viện xử lý PDF như PyMuPDF
\cite{pymupdf2026}. Với PDF được quét từ ảnh, nhận dạng ký tự quang học được
dùng để lấy văn bản và vị trí từ \cite{du2020ppocr}. Đầu ra của tầng này là các
block có trang, bbox, loại block, thứ tự đọc và ngữ cảnh đề mục, làm nền cho
chia đoạn, lập chỉ mục, truy xuất và dẫn chứng nguồn.

\section{Trích xuất và tái dựng bảng}
\label{section:table-extraction-foundation}

Bảng trong PDF không chỉ chứa văn bản mà còn chứa quan hệ hàng, cột và ô. Nếu
bảng bị nối thành văn bản phẳng, hệ thống có thể giữ được từ khóa nhưng mất
quan hệ điều kiện--giá trị. Vì vậy, đồ án xem trích xuất bảng là bước tái dựng cấu
trúc, không chỉ là lấy chữ trong vùng bảng.

Table Transformer được xây dựng trên hướng phát hiện đối tượng bằng
Transformer, kế thừa các ý tưởng từ Transformer và DETR
\cite{vaswani2017transformer,carion2020detr}. PubTables-1M cung cấp dữ liệu
và mô hình cho bài toán phát hiện bảng, nhận diện cấu trúc hàng, cột và ô
\cite{smock2022pubtables}. Trong đồ án, TATR được dùng để nhận diện cấu trúc
hình học của bảng, còn Hybrid TATR kết hợp cấu trúc đó với các hộp từ lấy từ lớp
văn bản PDF hoặc OCR để gán nội dung vào ô.

Trong triển khai hiện tại, Hybrid TATR trong quy trình chính ưu tiên dùng hộp từ
từ lớp văn bản PDF. Với tài liệu scan hoặc ảnh bảng, hộp từ OCR được xem là
nguồn đầu vào bổ sung trong thí nghiệm minh họa và hướng mở rộng. Nói cách
khác, OCR cung cấp chữ và tọa độ từ, còn TATR cung cấp cấu trúc hình học của
bảng.

\begin{table}[H]
\centering
\small
\renewcommand{\arraystretch}{1.15}
\caption{Vai trò của các thành phần trong xử lý bảng}
\label{tab:table-foundation-roles}
\begin{tabularx}{\textwidth}{
    |>{\raggedright\arraybackslash}p{0.25\textwidth}
    |>{\raggedright\arraybackslash}X
    |>{\raggedright\arraybackslash}p{0.28\textwidth}|
}
\hline
\textbf{Thành phần} & \textbf{Vai trò} & \textbf{Giới hạn} \\
\hline
Phát hiện vùng bảng &
Xác định bbox bảng để crop ảnh bảng hoặc lấy các từ nằm trong vùng bảng. &
Bbox lệch có thể làm sai các bước sau. \\
\hline
TATR &
Nhận diện hàng, cột, ô và vùng liên quan của bảng. &
Không tự gán văn bản vào từng ô. \\
\hline
Hộp từ PDF/OCR &
Cung cấp chữ và tọa độ từ để gán vào ô bảng. &
Phụ thuộc chất lượng lớp văn bản PDF hoặc OCR. \\
\hline
Hybrid TATR &
Kết hợp cấu trúc hình học với hộp từ để tạo bảng có nội dung. &
Chưa bảo đảm tái dựng hoàn hảo mọi bảng phức tạp. \\
\hline
\end{tabularx}
\end{table}

\section{Chia đoạn, siêu dữ liệu và lập chỉ mục}
\label{section:chunking-indexing-foundation}

Chia đoạn là bước chuyển block thành đơn vị truy xuất. Trong đồ án, đơn vị truy
xuất không chỉ là một đoạn cắt theo độ dài, mà còn mang siêu dữ liệu như tài liệu
nguồn, trang, loại block, đường dẫn đề mục, bbox và thông tin bảng nếu có. Các
siêu dữ liệu này giúp hệ thống truy vết bằng chứng và tạo dẫn chứng về đúng vị
trí trong PDF.

Với văn bản thông thường, đơn vị truy xuất cần đủ ngắn để truy xuất chính xác
nhưng vẫn giữ ngữ cảnh đề mục. Với bảng, hệ thống có thể sinh nhiều dạng đơn vị
như tóm tắt bảng, cấu trúc bảng, hàng bảng và ô bảng. Cách biểu diễn này giúp
truy xuất tìm đúng bảng, hàng hoặc ô thay vì chỉ tìm đúng trang chứa bảng.

Lập chỉ mục là bước xây dựng cấu trúc tìm kiếm từ các đơn vị truy xuất. Với truy xuất từ
khóa, chỉ mục phục vụ tính điểm theo từ. Với truy xuất ngữ nghĩa, mỗi đơn vị được
mã hóa thành biểu diễn vector. Các thư viện tìm kiếm vector như FAISS hỗ trợ tìm
kiếm tương đồng trên tập vector lớn \cite{johnson2021faiss}. Trong triển khai của
đồ án, chỉ mục vector được lưu bằng các biểu diễn đã chuẩn hóa; khi môi trường
hỗ trợ, hệ thống có thể ghi thêm tệp chỉ mục FAISS.

\section{Truy xuất thông tin}
\label{section:retrieval-foundation}

Truy xuất thông tin là bước tìm các đơn vị nội dung liên quan đến câu hỏi. Chất lượng
truy xuất quyết định chất lượng bằng chứng đầu vào cho tầng hỏi đáp. Đồ án so sánh
ba nhóm chính: truy xuất theo từ khóa, truy xuất ngữ nghĩa và truy xuất kết hợp.

\subsection{Truy xuất theo từ khóa}

BM25 là phương pháp xếp hạng theo từ khóa dựa trên tần suất từ, độ hiếm của từ
trong tập dữ liệu và độ dài tài liệu \cite{robertson2009bm25}. Với truy vấn \(q\) và
đơn vị truy xuất \(d\), điểm BM25 có thể viết như sau:

\begin{equation}
\label{eq:bm25}
\operatorname{BM25}(q,d) =
\sum_{t \in q}
\operatorname{IDF}(t)
\cdot
\frac{f(t,d)(k_1+1)}
{f(t,d)+k_1\left(1-b+b\frac{|d|}{\operatorname{avgdl}}\right)}
\end{equation}

\begin{equation}
\label{eq:idf}
\operatorname{IDF}(t)=
\log
\frac{N-n(t)+0.5}{n(t)+0.5}
\end{equation}

Trong đó \(f(t,d)\) là số lần từ \(t\) xuất hiện trong đơn vị truy xuất \(d\),
\(|d|\) là độ dài của đơn vị đó, \(\operatorname{avgdl}\) là độ dài trung bình của
các đơn vị truy xuất, \(N\) là số đơn vị trong tập dữ liệu và \(n(t)\) là số đơn vị
chứa từ \(t\). Trong đồ án, BM25 phù
hợp với câu hỏi có thuật ngữ rõ ràng như tên điều khoản, loại điểm, tín chỉ hoặc
cụm quy định.

\subsection{Truy xuất ngữ nghĩa}

Truy xuất ngữ nghĩa biểu diễn câu hỏi và đơn vị truy xuất trong cùng không gian vector.
Dense Passage Retrieval và Sentence-BERT là hai hướng tiêu biểu cho truy xuất
bằng biểu diễn vector \cite{karpukhin2020dpr,reimers2019sentencebert}. Gọi \(e_q\)
là vector của câu hỏi và \(e_d\) là vector của đơn vị truy xuất, độ tương đồng có
thể tính bằng cosine:

\begin{equation}
\label{eq:cosine-similarity}
\operatorname{sim}(q,d)=
\frac{e_q \cdot e_d}{\|e_q\|\|e_d\|}
\end{equation}

Khi các vector đã được chuẩn hóa, điểm dense có thể tính bằng tích vô hướng:

\begin{equation}
\label{eq:dense-dot-product}
\operatorname{score}_{dense}(q,d)=e_q \cdot e_d
\end{equation}

Truy xuất ngữ nghĩa hỗ trợ các câu hỏi diễn đạt tự nhiên, nhưng có thể lấy nhầm
đoạn gần nghĩa mà thiếu điều kiện hoặc giá trị chính xác. Vì vậy, đồ án không dùng
truy xuất ngữ nghĩa đơn lẻ làm cấu hình chính.

\subsection{Truy xuất kết hợp và sắp xếp lại}

Truy xuất kết hợp dùng cả tín hiệu từ khóa và tín hiệu ngữ nghĩa. Trong đồ án,
điểm hybrid được tính bằng cách chuẩn hóa và cộng trọng số:

\begin{equation}
\label{eq:hybrid-score}
\operatorname{score}_{hybrid}(q,d)
=
\alpha \cdot \widehat{\operatorname{score}}_{BM25}(q,d)
+
(1-\alpha) \cdot \widehat{\operatorname{score}}_{dense}(q,d)
\end{equation}

Ngoài ra, hệ thống hỗ trợ Reciprocal Rank Fusion để hợp nhất danh sách xếp hạng
\cite{cormack2009rrf}:

\begin{equation}
\label{eq:rrf}
\operatorname{RRF}(d)=
\sum_{r \in R}
\frac{1}{k+\operatorname{rank}_r(d)}
\end{equation}

Trong đó \(R\) là tập các danh sách xếp hạng và \(\operatorname{rank}_r(d)\) là vị trí
của đơn vị truy xuất \(d\) trong danh sách \(r\). Sau bước truy xuất ban đầu, hệ thống có thể
sắp xếp lại một tập ứng viên nhỏ. ColBERT minh họa hướng tương tác muộn giữa
truy vấn và tài liệu \cite{khattab2020colbert}. Trong mã nguồn của đồ án, bộ sắp
xếp lại có thể là luật kinh nghiệm, cross-encoder hoặc ColBERT tùy cấu hình:

\begin{equation}
\label{eq:rerank-score}
\operatorname{score}_{rerank}(q,d)=g_{\theta}(q,d)
\end{equation}

\section{Hỏi đáp có căn cứ và dẫn chứng nguồn}
\label{section:grounded-qa-foundation}

Hỏi đáp tăng cường bằng truy xuất sử dụng các bằng chứng được truy xuất làm ngữ
cảnh cho quá trình sinh câu trả lời \cite{lewis2020rag}. Cách tiếp cận này phù hợp
với PDF vì câu trả lời cần bám vào tài liệu nguồn thay vì chỉ dựa vào kiến thức có
sẵn của mô hình ngôn ngữ.

Trong đồ án, câu trả lời được xem là có căn cứ khi nội dung trả lời được hỗ trợ bởi
bằng chứng đã truy xuất và có dẫn chứng kèm theo. Quy trình vì vậy gồm truy xuất
top-k đơn vị nội dung, kiểm tra bằng chứng, sinh câu trả lời và tạo dẫn chứng. Các hướng như
Self-RAG và Adaptive-RAG cũng nhấn mạnh vai trò của kiểm soát truy xuất, sinh
và tự đánh giá câu trả lời \cite{asai2024selfrag,jeong2024adaptiverag}.

Trong triển khai hiện tại, định tuyến câu hỏi và kiểm tra bằng chứng được thực hiện
bằng luật và siêu dữ liệu, chưa huấn luyện mô hình phân loại truy vấn mới. Cách làm
này phù hợp với phạm vi đồ án vì dễ giải thích, dễ phân tích lỗi và thuận tiện cho
thí nghiệm loại bỏ thành phần.

\section{Độ đo đánh giá}
\label{section:evaluation-foundation}

Đồ án đánh giá hệ thống theo nhiều tầng: tiếp nhận và trích xuất PDF, trích xuất
bảng, truy xuất, hỏi đáp có căn cứ, dẫn chứng nguồn và hỏi đáp trên bảng. Cách
đánh giá theo tầng giúp phân biệt lỗi đến từ trích xuất tài liệu, truy xuất bằng chứng
hay sinh câu trả lời.

\subsection{Độ đo tiếp nhận, OCR và hiệu năng}

Với tầng tiếp nhận và trích xuất PDF, một độ đo cơ bản là tỷ lệ xử lý thành công trên tập mẫu. Gọi
\(M\) là số tài liệu hoặc số trang được đánh giá, và \(s_i \in \{0,1\}\) là biến
cho biết mẫu \(i\) có được xử lý thành công hay không. Khi đó:

\begin{equation}
\label{eq:success-rate}
\operatorname{SuccessRate} =
\frac{1}{M}
\sum_{i=1}^{M} s_i
\end{equation}

Đối với OCR, đồ án dùng Token F1 và CER. Sau khi chuẩn hóa văn bản, gọi
\(TP_{tok}\), \(FP_{tok}\) và \(FN_{tok}\) lần lượt là số token đúng, số token thừa
và số token thiếu. Khi đó:

\begin{equation}
\label{eq:token-precision-recall}
\operatorname{Precision}_{tok} =
\frac{TP_{tok}}{TP_{tok}+FP_{tok}},
\qquad
\operatorname{Recall}_{tok} =
\frac{TP_{tok}}{TP_{tok}+FN_{tok}}
\end{equation}

\begin{equation}
\label{eq:token-f1}
\operatorname{TokenF1} =
\frac{2 \cdot \operatorname{Precision}_{tok} \cdot \operatorname{Recall}_{tok}}
{\operatorname{Precision}_{tok}+\operatorname{Recall}_{tok}}
\end{equation}

Gọi \(c_i^{pred}\) là chuỗi ký tự OCR dự đoán, \(c_i^{gold}\) là chuỗi chuẩn và
\(\operatorname{ED}(\cdot,\cdot)\) là khoảng cách chỉnh sửa Levenshtein, CER được
tính bởi:

\begin{equation}
\label{eq:cer}
\operatorname{CER} =
\frac{\sum_{i=1}^{M} \operatorname{ED}(c_i^{pred}, c_i^{gold})}
{\sum_{i=1}^{M} |c_i^{gold}|}
\end{equation}

Đối với thời gian xử lý, đồ án báo cáo độ trễ trung bình trên tập đánh giá. Nếu
\(t_i\) là thời gian xử lý của mẫu \(i\), thì:

\begin{equation}
\label{eq:avg-latency}
\operatorname{Latency}_{avg} =
\frac{1}{M}
\sum_{i=1}^{M} t_i
\end{equation}

Success Rate cho biết quy trình có hoàn tất xử lý hay không, Token F1 và CER đo
chất lượng OCR ở mức nội dung chữ, còn \(\operatorname{Latency}_{avg}\) phản ánh
chi phí thực thi của cấu hình đang xét.

\subsection{Độ đo truy xuất}

Các độ đo truy xuất chính gồm Hit@k, Recall@k, MRR và NDCG. NDCG đo chất
lượng thứ hạng có xét vị trí của kết quả liên quan \cite{jarvelin2002ndcg}. BEIR,
SciFact và QASPER được dùng làm cơ sở dữ liệu benchmark liên quan đến truy
xuất và hỏi đáp trên dữ liệu khoa học
\cite{thakur2021beir,wadden2020scifact,dasigi2021qasper}.

Gọi \(Q\) là tập câu hỏi, \(G_i\) là tập bằng chứng đúng của câu hỏi \(i\), và
\(R_i@k\) là top-\(k\) kết quả hệ thống trả về:

\begin{equation}
\label{eq:hit-at-k}
\operatorname{Hit@k} =
\frac{1}{|Q|}
\sum_{i=1}^{|Q|}
\mathbf{1}\left(R_i@k \cap G_i \neq \emptyset\right)
\end{equation}

\begin{equation}
\label{eq:recall-at-k}
\operatorname{Recall@k} =
\frac{1}{|Q|}
\sum_{i=1}^{|Q|}
\frac{|R_i@k \cap G_i|}{|G_i|}
\end{equation}

\begin{equation}
\label{eq:mrr}
\operatorname{MRR} =
\frac{1}{|Q|}
\sum_{i=1}^{|Q|}
\frac{1}{\operatorname{rank}_i}
\end{equation}

\begin{equation}
\label{eq:ndcg}
\operatorname{DCG@k} =
\sum_{j=1}^{k}
\frac{rel_j}{\log_2(j+1)},
\qquad
\operatorname{NDCG@k} =
\frac{\operatorname{DCG@k}}{\operatorname{IDCG@k}}
\end{equation}

Hit@k cho biết trong \(k\) kết quả đầu tiên có ít nhất một bằng chứng đúng hay
không. Recall@k đo tỷ lệ bằng chứng đúng được thu hồi trong top-\(k\). MRR nhấn
mạnh vị trí của bằng chứng đúng đầu tiên, còn NDCG@k đánh giá chất lượng thứ hạng
khi các kết quả đúng xuất hiện ở các vị trí khác nhau. Trong công thức NDCG,
\(rel_j\) là mức liên quan của kết quả tại vị trí \(j\), còn \(\operatorname{IDCG@k}\)
là điểm DCG lý tưởng nếu các kết quả được xếp theo thứ tự tốt nhất.

\subsection{Độ đo hỏi đáp có căn cứ}

Với \(N\) câu hỏi, các độ đo QA trong đồ án được dùng ở dạng tổng quát:

\begin{equation}
\label{eq:answer-match}
\operatorname{AnswerMatch} =
\frac{1}{N}
\sum_{i=1}^{N}
\mathbf{1}\left(\operatorname{match}(a_i,\hat{a}_i)=1\right)
\end{equation}

\begin{equation}
\label{eq:evidence-match}
\operatorname{EvidenceMatch} =
\frac{1}{N}
\sum_{i=1}^{N}
\mathbf{1}\left(\hat{E}_i \cap E_i \neq \emptyset\right)
\end{equation}

\begin{equation}
\label{eq:grounded-rate}
\operatorname{GroundedRate} =
\frac{\#\{\text{câu trả lời được bằng chứng hỗ trợ}\}}{N}
\end{equation}

\begin{equation}
\label{eq:hallucination-rate}
\operatorname{HallucinationRate} =
\frac{\#\{\text{câu trả lời chứa thông tin thiếu căn cứ}\}}{N}
\end{equation}

Trong đó \(a_i\) là câu trả lời chuẩn, \(\hat{a}_i\) là câu trả lời của hệ thống,
\(E_i\) là tập bằng chứng đúng và \(\hat{E}_i\) là bằng chứng hệ thống truy xuất hoặc
trích dẫn.
Answer Match đánh giá nội dung câu trả lời cuối cùng, Evidence Match đánh giá
hệ thống có tìm hoặc trích dẫn đúng bằng chứng hay không. Grounded Rate đo tỷ lệ
câu trả lời được hỗ trợ bởi bằng chứng, còn Hallucination Rate đo tỷ lệ câu trả lời có
thông tin không được bằng chứng hỗ trợ. Các quy tắc khớp cụ thể theo từng tập đánh giá
được trình bày trong Chương 6.

\subsection{Độ đo xử lý bảng}

Với phát hiện vùng bảng hoặc ô bảng, độ chồng khớp giữa bbox dự đoán \(B_p\)
và bbox ground truth \(B_g\) được đo bằng IoU:

\begin{equation}
\label{eq:iou}
\operatorname{IoU}(B_p,B_g)=
\frac{|B_p \cap B_g|}{|B_p \cup B_g|}
\end{equation}

Từ các cặp dự đoán--ground truth được ghép đúng theo ngưỡng IoU, hệ thống tính
Precision, Recall và F1:

\begin{equation}
\label{eq:table-precision-recall}
\operatorname{Precision} =
\frac{TP}{TP+FP},
\qquad
\operatorname{Recall} =
\frac{TP}{TP+FN}
\end{equation}

\begin{equation}
\label{eq:table-f1}
\operatorname{F1} =
\frac{2 \cdot \operatorname{Precision} \cdot \operatorname{Recall}}
{\operatorname{Precision}+\operatorname{Recall}}
\end{equation}

Precision cho biết trong các dự đoán của hệ thống có bao nhiêu dự đoán đúng,
Recall cho biết hệ thống tìm được bao nhiêu đối tượng đúng trong ground truth,
còn F1 là trung bình điều hòa giữa hai độ đo này. Khi đánh giá phát hiện bảng,
một dự đoán chỉ được tính là đúng nếu bbox dự đoán khớp với bbox ground truth
theo ngưỡng IoU đã chọn.

Với tái dựng bảng, Exact CSV kiểm tra mức khớp toàn bộ bảng sau chuẩn hóa:

\begin{equation}
\label{eq:exact-csv}
\operatorname{ExactCSV} =
\frac{1}{M}
\sum_{i=1}^{M}
\mathbf{1}
\left(
\operatorname{norm}(CSV_i^{pred}) =
\operatorname{norm}(CSV_i^{gold})
\right)
\end{equation}

Ngoài Exact CSV, đồ án còn dùng thêm Cell F1, Structure F1 và Text Assignment
F1 để phân tích chi tiết hơn. Với Cell F1, gọi \(C^{pred}\) và \(C^{gold}\) là tập
ô dự đoán và ô chuẩn sau khi ghép đúng theo ngưỡng IoU \(\tau\). Khi đó:

\begin{equation}
\label{eq:cell-f1}
\operatorname{CellF1@}\tau =
\frac{2 \cdot \operatorname{Precision}_{cell} \cdot \operatorname{Recall}_{cell}}
{\operatorname{Precision}_{cell}+\operatorname{Recall}_{cell}}
\end{equation}

trong đó:

\begin{equation}
\label{eq:cell-precision-recall}
\operatorname{Precision}_{cell} =
\frac{|C^{pred} \cap C^{gold}|}{|C^{pred}|},
\qquad
\operatorname{Recall}_{cell} =
\frac{|C^{pred} \cap C^{gold}|}{|C^{gold}|}
\end{equation}

Structure F1 đo mức khớp của các quan hệ cấu trúc trong bảng, chẳng hạn hàng,
cột, span hoặc quan hệ kề nhau giữa các ô. Gọi \(S^{pred}\) và \(S^{gold}\) là tập
quan hệ cấu trúc dự đoán và chuẩn, khi đó:

\begin{equation}
\label{eq:structure-f1}
\operatorname{StructureF1} =
\frac{2 \cdot \operatorname{Precision}_{struct} \cdot \operatorname{Recall}_{struct}}
{\operatorname{Precision}_{struct}+\operatorname{Recall}_{struct}}
\end{equation}

\begin{equation}
\label{eq:structure-precision-recall}
\operatorname{Precision}_{struct} =
\frac{|S^{pred} \cap S^{gold}|}{|S^{pred}|},
\qquad
\operatorname{Recall}_{struct} =
\frac{|S^{pred} \cap S^{gold}|}{|S^{gold}|}
\end{equation}

Text Assignment F1 đo mức khớp giữa văn bản và ô được gán. Gọi \(A^{pred}\) và
\(A^{gold}\) là tập các cặp gán \((\text{token}, \text{cell})\) của hệ thống và
ground truth, ta có:

\begin{equation}
\label{eq:text-assignment-f1}
\operatorname{TextAssignmentF1} =
\frac{2 \cdot \operatorname{Precision}_{assign} \cdot \operatorname{Recall}_{assign}}
{\operatorname{Precision}_{assign}+\operatorname{Recall}_{assign}}
\end{equation}

\begin{equation}
\label{eq:text-assignment-precision-recall}
\operatorname{Precision}_{assign} =
\frac{|A^{pred} \cap A^{gold}|}{|A^{pred}|},
\qquad
\operatorname{Recall}_{assign} =
\frac{|A^{pred} \cap A^{gold}|}{|A^{gold}|}
\end{equation}

Exact CSV là độ đo nghiêm ngặt nhất vì chỉ một sai khác nhỏ về cấu trúc hoặc thứ
tự ô cũng làm toàn bộ bảng bị tính sai. Cell F1, Structure F1 và
Text Assignment F1 giúp phân biệt lỗi nằm ở phát hiện ô, ở quan hệ cấu trúc hay
ở bước gán văn bản vào ô.

\subsection{Độ đo truy xuất bảng}

Với hỏi đáp trên bảng, đồ án đánh giá liệu top-\(k\) kết quả có chứa đúng đơn vị cần truy
xuất ở mức bảng, hàng, cột hoặc ô hay không. Gọi \(u \in
\{\text{table},\text{row},\text{column},\text{cell}\}\) là loại đơn vị cần đánh
giá, \(G_i^{(u)}\) là tập đơn vị đúng của câu hỏi \(i\), và \(R_i^{(u)}@k\) là tập
đơn vị loại \(u\) xuất hiện trong top-\(k\) kết quả. Khi đó:

\begin{equation}
\label{eq:unit-match-at-k}
\operatorname{Match}_{u}@k =
\frac{1}{N}
\sum_{i=1}^{N}
\mathbf{1}\left(R_i^{(u)}@k \cap G_i^{(u)} \neq \emptyset\right)
\end{equation}

Từ công thức tổng quát này, đồ án suy ra các độ đo Table Match@k, Row Match@k,
Column Match@k và Cell Match@k bằng cách lần lượt thay \(u\) bởi
\(\text{table}\), \(\text{row}\), \(\text{column}\) và \(\text{cell}\). Các độ đo
này cho biết hệ thống đã tìm đúng mức bằng chứng nào trong top-\(k\): đúng bảng
nhưng sai hàng, đúng hàng nhưng sai ô, hoặc đúng trực tiếp ô cần tra cứu.

\section{Công nghệ sử dụng trong đồ án}
\label{section:technology-foundation}

Hệ thống được triển khai bằng Python và chia thành các mô-đun tương ứng với các
tầng của quy trình. Bảng dưới tóm tắt nhóm công nghệ và nơi liên hệ trong mã nguồn.

\begin{table}[H]
\centering
\small
\renewcommand{\arraystretch}{1.15}
\caption{Nhóm công nghệ và vai trò trong đồ án}
\label{tab:chapter3-technology-summary}
\begin{tabularx}{\textwidth}{
    |>{\raggedright\arraybackslash}p{0.25\textwidth}
    |>{\raggedright\arraybackslash}X
    |>{\raggedright\arraybackslash}p{0.27\textwidth}|
}
\hline
\textbf{Nhóm công nghệ} & \textbf{Vai trò} & \textbf{Liên hệ triển khai} \\
\hline
Xử lý PDF &
Đọc lớp văn bản, lấy tọa độ, kết xuất trang và chuẩn hóa block. &
\texttt{app/ingest}, \texttt{app/loaders} \\
\hline
OCR &
Nhận dạng văn bản và hộp từ trong tài liệu ảnh hoặc vùng cần nhận dạng. &
\texttt{app/ingest/extract/ocr.py} \\
\hline
Xử lý bảng &
Phát hiện vùng bảng, tái dựng cấu trúc và gán văn bản vào ô. &
\texttt{app/ingest/extract/table.py}, \texttt{hybrid\_tatr\_table.py} \\
\hline
Truy xuất thông tin &
BM25, truy xuất ngữ nghĩa, truy xuất kết hợp và sắp xếp lại kết quả. &
\texttt{app/retrieval} \\
\hline
Hỏi đáp có căn cứ &
Định tuyến câu hỏi, kiểm tra bằng chứng, sinh câu trả lời và tạo dẫn chứng. &
\texttt{app/qa} \\
\hline
Đánh giá và chạy thử &
Chạy đánh giá từng tầng và hỏi đáp nhanh qua giao diện dòng lệnh. &
\texttt{scripts/benchmark\_*.py}, \texttt{quick\_qa\_terminal.py} \\
\hline
\end{tabularx}
\end{table}

\section{Kết chương}

Chương này đã tóm tắt các nền tảng được dùng trực tiếp trong đồ án: xử lý PDF,
trích xuất bảng, chia đoạn, lập chỉ mục, truy xuất thông tin, hỏi đáp có căn cứ và
các độ đo đánh giá. Các nội dung được trình bày ở mức vừa đủ để làm cơ sở cho
Chương 4 về luồng tiếp nhận và trích xuất PDF, Chương 5 về luồng xử lý truy vấn
và Chương 6 về thực nghiệm.

\end{document}
